\shead{Sơ đồ Venn hai vòng tròn}
% ---- chapter 12 helpers: Venn diagrams with two and three circles ----
\newcommand{\stpVennII}[6]{% #1 #2 circle labels; #3 only A, #4 both, #5 only B, #6 outside (text)
\begin{tikzpicture}[font=\small]
\fill[tmBlue!14] (0,0) circle (1.3);
\fill[tmOrange!18] (1.7,0) circle (1.3);
\begin{scope}\clip (0,0) circle (1.3);\fill[tmGreen!35] (1.7,0) circle (1.3);\end{scope}
\draw[line width=0.9pt,tmBlue!80!black] (0,0) circle (1.3);
\draw[line width=0.9pt,tmOrange!85!black] (1.7,0) circle (1.3);
\draw[line width=0.8pt,Slate,rounded corners=4pt] (-2.0,-1.75) rectangle (3.7,1.95);
\node[font=\footnotesize\bfseries,text=tmBlue!80!black] at (-0.55,1.62) {#1};
\node[font=\footnotesize\bfseries,text=tmOrange!85!black] at (2.25,1.62) {#2};
\node at (-0.45,0) {#3}; \node at (0.85,0) {#4}; \node at (2.15,0) {#5};
\node[anchor=south east,font=\footnotesize] at (3.6,-1.7) {#6};
\end{tikzpicture}}
\newcommand{\stpVennIII}[9]{% #1 #2 #3 labels A B C; #4 only A, #5 only B, #6 only C, #7 A&B only, #8 A&C only, #9 B&C only
\def\stpVa{#1}\def\stpVb{#2}\def\stpVc{#3}\def\stpVd{#4}\def\stpVe{#5}\def\stpVf{#6}\def\stpVg{#7}\def\stpVh{#8}\def\stpVi{#9}\stpVennIIIx}
\newcommand{\stpVennIIIx}[2]{% #1 all three, #2 outside (text)
\begin{tikzpicture}[font=\small]
\fill[tmBlue!30,fill opacity=0.55] (-0.7,0.4) circle (1.15);
\fill[tmOrange!35,fill opacity=0.55] (0.7,0.4) circle (1.15);
\fill[tmGreen!30,fill opacity=0.55] (0,-0.8) circle (1.15);
\draw[line width=0.9pt,tmBlue!80!black] (-0.7,0.4) circle (1.15);
\draw[line width=0.9pt,tmOrange!85!black] (0.7,0.4) circle (1.15);
\draw[line width=0.9pt,tmGreen!70!black] (0,-0.8) circle (1.15);
\draw[line width=0.8pt,Slate,rounded corners=4pt] (-4.0,-2.3) rectangle (4.0,1.85);
\node[font=\footnotesize\bfseries,text=tmBlue!80!black,anchor=east] at (-1.85,1.0) {\stpVa};
\node[font=\footnotesize\bfseries,text=tmOrange!85!black,anchor=west] at (1.85,1.0) {\stpVb};
\node[font=\footnotesize\bfseries,text=tmGreen!60!black,anchor=west] at (1.2,-1.65) {\stpVc};
\node at (-1.25,0.65) {\stpVd}; \node at (1.25,0.65) {\stpVe}; \node at (0,-1.45) {\stpVf};
\node at (0,0.98) {\stpVg}; \node at (-0.72,-0.45) {\stpVh}; \node at (0.72,-0.45) {\stpVi};
\node at (0,0.05) {#1};
\node[anchor=south west,font=\footnotesize] at (-3.9,-2.25) {#2};
\end{tikzpicture}}
\begin{knowbox}
\textbf{Sơ đồ Venn} biểu diễn các nhóm bằng những vòng khép kín (vòng tròn) nằm trong một hình chữ nhật biểu thị toàn bộ.
\begin{center}
\stpVennII{Nhóm A}{Nhóm B}{\footnotesize chỉ A}{\footnotesize cả hai}{\footnotesize chỉ B}{không thuộc nhóm nào}
\end{center}
\begin{itemize}[leftmargin=1.4em,itemsep=3pt]
\item \emph{Phần giao nhau} của hai vòng tròn gồm những người thuộc \textbf{cả hai} nhóm.
\item Phần còn lại của mỗi vòng tròn (phần khác biệt) gồm những người \textbf{chỉ} thuộc nhóm đó.
\item Phần nằm trong hình chữ nhật nhưng ở ngoài hai vòng tròn gồm những người \textbf{không thuộc nhóm nào}.
\end{itemize}
Khi cộng A $+$ B, những người thuộc cả hai nhóm được tính \emph{hai lần}. Vì vậy:
\begin{center}
\textbf{thuộc ít nhất một nhóm} $=$ A $+$ B $-$ cả hai \qquad\quad \textbf{cả hai} $=$ A $+$ B $-$ (thuộc ít nhất một nhóm)\\[3pt]
\textbf{chỉ A} $=$ A $-$ cả hai \qquad\quad \textbf{không thuộc nhóm nào} $=$ toàn bộ $-$ (thuộc ít nhất một nhóm)
\end{center}
\end{knowbox}
\shead{Sơ đồ Venn ba vòng tròn}
\begin{knowbox}
Ba vòng tròn A, B, C tạo thành $7$ phần ở bên trong (cùng với phần ở bên ngoài). Ta điền số vào sơ đồ \textbf{từ trong ra ngoài}:
\begin{enumerate}[leftmargin=1.6em,itemsep=2pt]
\item phần chính giữa: những người thuộc cả ba nhóm;
\item mỗi phần giao của hai vòng: (thuộc cả hai nhóm đó) $-$ (thuộc cả ba nhóm);
\item mỗi phần ``chỉ'': (cả vòng tròn) $-$ (ba số đã điền trong vòng tròn đó).
\end{enumerate}
Cộng tất cả các phần ta được số người thuộc ít nhất một nhóm. Viết gọn:
\begin{center}
thuộc ít nhất một nhóm $=$ A $+$ B $+$ C $-$ (A và B) $-$ (B và C) $-$ (C và A) $+$ (cả ba).
\end{center}
\end{knowbox}
\shead{Tính ngược bằng cách lập bảng}
\begin{knowbox}
Khi đồ vật được chuyển từ người này (hộp này) sang người khác nhiều lần và ta biết kết quả \textbf{cuối cùng}, ta \textbf{tính ngược} từ cuối:
\begin{itemize}[leftmargin=1.4em,itemsep=3pt]
\item Chuyển cho nhau \emph{không} làm thay đổi tổng số. Nếu lúc cuối mọi người có số lượng bằng nhau thì mỗi người có: tổng số $:$ số người.
\item Xét các lần chuyển theo thứ tự ngược lại: người \emph{đã cho} được nhận lại ($+$), người \emph{đã nhận} phải trả lại ($-$).
\item ``Cho bạn đúng bằng số bạn đang có'' làm số của người nhận \emph{gấp đôi}; khi tính ngược ta \emph{chia đôi}.
\item Lập \textbf{bảng}, mỗi hàng ứng với một thời điểm, để trình bày gọn gàng. Thử lại bằng cách tính xuôi.
\end{itemize}
\end{knowbox}
\gap
\begin{exbox}{\ExWord\ 1}
Một nhóm $30$ học sinh, mỗi bạn tham gia ít nhất một trong hai câu lạc bộ: $18$ bạn tham gia CLB Cờ vua và $19$ bạn tham gia CLB Vẽ. Hỏi có bao nhiêu bạn tham gia cả hai câu lạc bộ? Bao nhiêu bạn chỉ tham gia một câu lạc bộ?
\guide
Khi lấy $18+19=37$, những bạn tham gia cả hai câu lạc bộ được tính hai lần, nên
\par số bạn tham gia cả hai câu lạc bộ là: $37-30=7$ (bạn).\par
Chỉ tham gia Cờ vua: $18-7=11$ (bạn); chỉ tham gia Vẽ: $19-7=12$ (bạn); chỉ tham gia một câu lạc bộ: $11+12=23$ (bạn).
\begin{center}
\stpVennII{Cờ vua $18$}{Vẽ $19$}{$11$}{$\mathbf{7}$}{$12$}{$30$ học sinh}
\end{center}
\dapso{$7$ bạn; $23$ bạn.}
\end{exbox}
\gap
\begin{exbox}{\ExWord\ 2}
Lớp 5C có $36$ học sinh. Có $20$ bạn biết bơi, $15$ bạn biết chơi cầu lông và $8$ bạn biết cả hai. Hỏi có bao nhiêu bạn không biết bơi cũng không biết chơi cầu lông?
\guide
\emph{Cách 1.} Số bạn biết ít nhất một môn: $20+15-8=27$ (bạn).\par Không biết môn nào: $36-27=9$ (bạn).\par
\emph{Cách 2.} Chỉ biết bơi: $20-8=12$ (bạn); chỉ biết cầu lông: $15-8=7$ (bạn).\par Biết ít nhất một môn: $12+7+8=27$ (bạn); không biết môn nào: $36-27=9$ (bạn).
\begin{center}
\stpVennII{Bơi $20$}{Cầu lông $15$}{$12$}{$8$}{$7$}{không môn nào: $\mathbf{9}$}
\end{center}
\dapso{$9$ bạn.}
\end{exbox}
\gap
\begin{exbox}{\ExWord\ 3}
Mỗi bạn trong một nhóm học sinh thích ít nhất một trong ba môn: bóng đá, bơi, cờ vua. Có $25$ bạn thích bóng đá, $20$ bạn thích bơi, $18$ bạn thích cờ vua; $8$ bạn thích bóng đá và bơi, $6$ bạn thích bơi và cờ vua, $7$ bạn thích bóng đá và cờ vua, $3$ bạn thích cả ba môn.\par
\pa{a} Có bao nhiêu bạn chỉ thích bóng đá? \quad \pa{b} Nhóm có bao nhiêu học sinh?
\guide
Điền từ trong ra ngoài: cả ba môn: $3$ bạn.\par
Chỉ thích bóng đá và bơi: $8-3=5$; \ chỉ thích bơi và cờ vua: $6-3=3$; \ chỉ thích bóng đá và cờ vua: $7-3=4$.\par
\pa{a} Chỉ thích bóng đá: $25-5-4-3=13$ (bạn).\par
Chỉ thích bơi: $20-5-3-3=9$; \ chỉ thích cờ vua: $18-4-3-3=8$.
\begin{center}
\stpVennIII{Bóng đá $25$}{Bơi $20$}{Cờ vua $18$}{$\mathbf{13}$}{$9$}{$8$}{$5$}{$4$}{$3$}{$3$}{}
\end{center}
\pa{b} $13+9+8+5+4+3+3=45$ (học sinh).\par
(Thử lại: $25+20+18-8-6-7+3=45$.)
\dapso{\pa{a} $13$ bạn. \pa{b} $45$ học sinh.}
\end{exbox}
\gap
\begin{exbox}{\ExWord\ 4}
An, Bình và Chi có tất cả $75$ hình dán. An cho Bình $6$ hình, sau đó Bình cho Chi $9$ hình, rồi Chi cho An $4$ hình thì lúc này số hình dán của ba bạn bằng nhau. Hỏi lúc đầu mỗi bạn có bao nhiêu hình dán?
\guide
Tổng số vẫn là $75$, nên lúc cuối mỗi bạn có $75:3=25$ (hình). Tính ngược, mỗi hàng ứng với một lần cho:
\begin{center}
\renewcommand{\arraystretch}{1.25}
\begin{tabular}{|l|c|c|c|}
\hline
 & An & Bình & Chi\\\hline
Lúc cuối & $25$ & $25$ & $25$\\\hline
Trước khi Chi cho An $4$ hình & $25-4=21$ & $25$ & $25+4=29$\\\hline
Trước khi Bình cho Chi $9$ hình & $21$ & $25+9=34$ & $29-9=20$\\\hline
Trước khi An cho Bình $6$ hình (lúc đầu) & $21+6=\mathbf{27}$ & $34-6=\mathbf{28}$ & $\mathbf{20}$\\\hline
\end{tabular}
\end{center}
Thử lại: $27+28+20=75$.
\dapso{An $27$ hình, Bình $28$ hình, Chi $20$ hình dán.}
\end{exbox}
\gap
\begin{exbox}{\ExWord\ 5}
Lan, Mai và Cúc có $120$ vỏ ốc. Đầu tiên Lan cho Mai và Cúc mỗi bạn số vỏ ốc đúng bằng số vỏ ốc bạn đó đang có. Sau đó Mai cũng làm như vậy với Lan và Cúc, cuối cùng Cúc cũng làm như vậy với Lan và Mai. Lúc này số vỏ ốc của ba bạn bằng nhau. Hỏi lúc đầu mỗi bạn có bao nhiêu vỏ ốc?
\guide
Lúc cuối mỗi bạn có $120:3=40$ (vỏ ốc). Mỗi lần cho, số vỏ ốc của hai bạn được nhận \emph{gấp đôi}; khi tính ngược ta chia đôi số của hai bạn đó, còn bạn đã cho thì được nhận lại phần đã cho.
\begin{center}
\renewcommand{\arraystretch}{1.25}
\begin{tabular}{|l|c|c|c|}
\hline
 & Lan & Mai & Cúc\\\hline
Lúc cuối & $40$ & $40$ & $40$\\\hline
Trước khi Cúc cho & $20$ & $20$ & $40+20+20=80$\\\hline
Trước khi Mai cho & $10$ & $20+10+40=70$ & $40$\\\hline
Lúc đầu (trước khi Lan cho) & $10+35+20=\mathbf{65}$ & $\mathbf{35}$ & $\mathbf{20}$\\\hline
\end{tabular}
\end{center}
Thử lại (tính xuôi): $(65;\ 35;\ 20)\to(10;\ 70;\ 40)\to(20;\ 20;\ 80)\to(40;\ 40;\ 40)$.
\dapso{Lan $65$, Mai $35$, Cúc $20$ vỏ ốc.}
\end{exbox}
\dividerbar{Luyện tập}
\begin{probbox}{1}{\relax}
Mỗi học sinh lớp 5A đều được ít nhất một điểm $10$ môn Toán hoặc môn Tiếng Việt. Có $21$ bạn được điểm $10$ môn Toán, $17$ bạn được điểm $10$ môn Tiếng Việt và $6$ bạn được điểm $10$ cả hai môn. Hỏi lớp 5A có bao nhiêu học sinh?
\end{probbox}
\probsol{
$21+17-6=32$ (học sinh).
\dapso{$32$ học sinh.}}
\gap
\begin{probbox}{2}{\relax}
Mỗi bạn trong $40$ học sinh thích đọc sách hoặc thích vẽ (hoặc thích cả hai). Có $28$ bạn thích đọc sách và $22$ bạn thích vẽ. Hỏi có bao nhiêu bạn thích cả hai?
\end{probbox}
\probsol{
Tổng $28+22=50$ tính hai lần những bạn thích cả hai.\par Thích cả hai: $50-40=10$ (bạn).
\dapso{$10$ bạn.}}
\gap
\begin{probbox}{3}{\relax}
Trong một đoàn $45$ khách du lịch, có $30$ người nói được tiếng Anh, $20$ người nói được tiếng Pháp và $5$ người không nói được thứ tiếng nào trong hai thứ tiếng đó.\par
\pa{a} Có bao nhiêu người nói được cả hai thứ tiếng? \quad \pa{b} Bao nhiêu người chỉ nói được tiếng Anh? \quad \pa{c} chỉ nói được tiếng Pháp?
\end{probbox}
\probsol{
Nói được ít nhất một thứ tiếng: $45-5=40$ (người).\par
\pa{a} Cả hai thứ tiếng: $30+20-40=10$ (người).\par
\pa{b} Chỉ tiếng Anh: $30-10=20$ (người).\par
\pa{c} Chỉ tiếng Pháp: $20-10=10$ (người).
\dapso{\pa{a} $10$ người. \pa{b} $20$ người. \pa{c} $10$ người.}}
\gap
\needspace{7.5cm}
\begin{probbox}{4}{\relax}
Sơ đồ Venn cho biết đồ uống của những người dự một bữa tiệc.
\begin{center}
\stpVennII{Trà}{Cà phê}{$14$}{$6$}{$9$}{không uống: $3$}
\end{center}
\pa{a} Có bao nhiêu người uống trà? \quad \pa{b} Bao nhiêu người uống cà phê? \quad \pa{c} Bữa tiệc có bao nhiêu người?
\end{probbox}
\probsol{
\pa{a} Uống trà: $14+6=20$ (người).\par
\pa{b} Uống cà phê: $6+9=15$ (người).\par
\pa{c} $14+6+9+3=32$ (người).
\dapso{\pa{a} $20$ người. \pa{b} $15$ người. \pa{c} $32$ người.}}
\gap
\needspace{7.5cm}
\begin{probbox}{5}{\relax}
Mỗi thành viên của một đội thể thao chơi bóng đá hoặc bóng rổ (hoặc cả hai), như sơ đồ dưới đây.
\begin{center}
\stpVennII{Bóng đá}{Bóng rổ}{$12$}{$4$}{$9$}{}
\end{center}
\pa{a} Đội có bao nhiêu thành viên? \quad \pa{b} Có bao nhiêu bạn chơi bóng rổ?
\end{probbox}
\probsol{
\pa{a} $12+4+9=25$ (thành viên).\par
\pa{b} $4+9=13$ (bạn).
\dapso{\pa{a} $25$ thành viên. \pa{b} $13$ bạn.}}
\gap
\begin{probbox}{6}{\relax}
Đội tuyển Toán gồm $20$ học sinh được giao hai bài toán khó; bạn nào cũng làm được ít nhất một bài. Có $3$ bạn không làm được bài 1 và $5$ bạn không làm được bài 2. Hỏi có bao nhiêu bạn làm được cả hai bài?
\end{probbox}
\probsol{
$3$ bạn không làm được bài 1 là những bạn chỉ làm được bài 2.\par $5$ bạn không làm được bài 2 là những bạn chỉ làm được bài 1.\par
Làm được cả hai bài: $20-3-5=12$ (bạn).
\dapso{$12$ bạn.}}
\gap
\begin{probbox}{7}{\relax}
Lớp 5D có $38$ học sinh. Có $15$ bạn đăng kí học bơi, $20$ bạn đăng kí học bóng đá và $7$ bạn đăng kí cả hai môn. Hỏi có bao nhiêu bạn không đăng kí môn nào?
\end{probbox}
\probsol{
Đăng kí ít nhất một môn: $15+20-7=28$ (bạn).\par Không đăng kí môn nào: $38-28=10$ (bạn).
\dapso{$10$ bạn.}}
\gap
\begin{probbox}{8}{\relax}
Một lớp có $35$ học sinh, trong đó $14$ bạn nuôi chó, $11$ bạn nuôi mèo và $5$ bạn nuôi cả chó và mèo. Hỏi có bao nhiêu bạn không nuôi chó cũng không nuôi mèo?
\end{probbox}
\probsol{
Nuôi ít nhất một con vật: $14+11-5=20$ (bạn).\par Không nuôi con nào: $35-20=15$ (bạn).
\dapso{$15$ bạn.}}
\gap
\begin{probbox}{9}{\relax}
Khảo sát $60$ gia đình trong một thôn: $42$ gia đình có ti vi, $35$ gia đình có máy tính và gia đình nào cũng có ít nhất một trong hai thứ đó.\par
\pa{a} Có bao nhiêu gia đình có cả hai thứ? \quad \pa{b} chỉ có ti vi? \quad \pa{c} chỉ có máy tính?
\end{probbox}
\probsol{
\pa{a} $42+35-60=17$ (gia đình).\par
\pa{b} $42-17=25$ (gia đình).\par
\pa{c} $35-17=18$ (gia đình).
\dapso{\pa{a} $17$. \pa{b} $25$. \pa{c} $18$ gia đình.}}
\gap
\begin{probbox}{10}{\relax}
Trong các số từ $1$ đến $30$, có bao nhiêu số chia hết cho $2$ hoặc chia hết cho $3$?
\end{probbox}
\probsol{
Chia hết cho $2$: $30:2=15$ (số). Chia hết cho $3$: $30:3=10$ (số).\par
Chia hết cho cả $2$ và $3$, tức là chia hết cho $6$: $6;\ 12;\ 18;\ 24;\ 30$: $5$ số.
\begin{center}
\stpVennII{Chia hết cho $2$}{Chia hết cho $3$}{$10$}{$5$}{$5$}{}
\end{center}
Chia hết cho $2$ hoặc cho $3$: $15+10-5=20$ (số).
\dapso{$20$ số.}}
\gap
\begin{probbox}{11}{\relax}
Trong các số từ $1$ đến $50$, có bao nhiêu số không chia hết cho $2$ và cũng không chia hết cho $5$?
\end{probbox}
\probsol{
Chia hết cho $2$: $25$ số; chia hết cho $5$: $10$ số; chia hết cho cả hai (cho $10$): $5$ số.\par
Chia hết cho $2$ hoặc $5$: $25+10-5=30$ (số).\par Không chia hết cho số nào: $50-30=20$ (số).
\dapso{$20$ số.}}
\gap
\begin{probbox}{12}{\relax}
Trong một nhóm nhạc, $16$ bạn chơi đàn piano, $10$ bạn chơi đàn ghi-ta và $22$ bạn chơi ít nhất một trong hai loại đàn đó. Hỏi có bao nhiêu bạn chơi cả hai loại đàn? Bao nhiêu bạn chỉ chơi piano?
\end{probbox}
\probsol{
Chơi cả hai loại đàn: $16+10-22=4$ (bạn).\par Chỉ chơi piano: $16-4=12$ (bạn).
\dapso{$4$ bạn; $12$ bạn.}}
\gap
\begin{probbox}{13}{\relax}
Một lớp có $18$ bạn dự thi hát, $13$ bạn dự thi vẽ và $25$ bạn dự thi ít nhất một cuộc thi. Hỏi có bao nhiêu bạn chỉ dự thi đúng một cuộc thi?
\end{probbox}
\probsol{
Dự thi cả hai: $18+13-25=6$ (bạn).\par Chỉ dự thi một cuộc thi: $25-6=19$ (bạn).
\dapso{$19$ bạn.}}
\gap
\begin{probbox}{14}{\relax}
Hỏi $100$ học sinh về bữa sáng: $58$ bạn thích phở, $47$ bạn thích bánh mì và $15$ bạn không thích món nào trong hai món đó. Hỏi có bao nhiêu bạn thích cả hai món?
\end{probbox}
\probsol{
Thích ít nhất một món: $100-15=85$ (bạn).\par Thích cả hai món: $58+47-85=20$ (bạn).
\dapso{$20$ bạn.}}
\gap
\needspace{7.5cm}
\begin{probbox}{15}{\relax}
Sơ đồ cho biết số học sinh của một lớp tham gia các câu lạc bộ Toán, Tiếng Anh và Mĩ thuật.
\begin{center}
\stpVennIII{Toán}{Tiếng Anh}{Mĩ thuật}{$8$}{$7$}{$6$}{$3$}{$2$}{$4$}{$1$}{không CLB nào: $5$}
\end{center}
\pa{a} Câu lạc bộ Toán có bao nhiêu học sinh? \quad \pa{b} Có bao nhiêu bạn tham gia đúng một câu lạc bộ?\par
\pa{c} Lớp có bao nhiêu học sinh?
\end{probbox}
\probsol{
\pa{a} Cộng mọi số nằm trong vòng Toán: $8+3+2+1=14$ (học sinh).\par
\pa{b} $8+7+6=21$ (bạn).\par
\pa{c} $8+7+6+3+2+4+1+5=36$ (học sinh).
\dapso{\pa{a} $14$. \pa{b} $21$. \pa{c} $36$ học sinh.}}
\gap
\begin{probbox}{16}{\relax}
Mỗi học sinh trong đội dự thi của trường thi ít nhất một trong ba môn Toán, Tiếng Anh, Mĩ thuật: $25$ bạn thi Toán, $20$ bạn thi Tiếng Anh, $15$ bạn thi Mĩ thuật. Có $6$ bạn thi Toán và Tiếng Anh, $4$ bạn thi Tiếng Anh và Mĩ thuật, $5$ bạn thi Toán và Mĩ thuật, và không có bạn nào thi cả ba môn. Hỏi đội có bao nhiêu học sinh?
\end{probbox}
\probsol{
Tổng $25+20+15=60$ tính hai lần mỗi bạn thi hai môn.\par
Số học sinh của đội: $60-6-4-5=45$ (học sinh).
\dapso{$45$ học sinh.}}
\gap
\begin{probbox}{17}{\relax}
Hỏi $40$ học sinh về loại quả yêu thích: $20$ bạn thích táo, $18$ bạn thích chuối, $16$ bạn thích cam; $7$ bạn thích táo và chuối, $6$ bạn thích chuối và cam, $5$ bạn thích táo và cam, $2$ bạn thích cả ba loại. Hỏi có bao nhiêu bạn không thích loại nào trong ba loại quả đó?
\end{probbox}
\probsol{
Thích ít nhất một loại: $20+18+16-7-6-5+2=38$ (bạn).\par Không thích loại nào: $40-38=2$ (bạn).
\dapso{$2$ bạn.}}
\gap
\begin{probbox}{18}{\relax}
Một lớp có $30$ học sinh: $15$ bạn thích hát, $12$ bạn thích múa, $10$ bạn thích vẽ; $5$ bạn thích hát và múa, $4$ bạn thích múa và vẽ, $3$ bạn thích hát và vẽ, $1$ bạn thích cả ba.\par
\pa{a} Có bao nhiêu bạn chỉ thích hát? \quad \pa{b} Có bao nhiêu bạn không thích hoạt động nào?
\end{probbox}
\probsol{
\pa{a} Chỉ thích hát và múa: $5-1=4$; chỉ thích hát và vẽ: $3-1=2$.\par Chỉ thích hát: $15-4-2-1=8$ (bạn).\par
\pa{b} Thích ít nhất một hoạt động: $15+12+10-5-4-3+1=26$; không thích hoạt động nào: $30-26=4$ (bạn).
\dapso{\pa{a} $8$ bạn. \pa{b} $4$ bạn.}}
\gap
\begin{probbox}{19}{\relax}
Trong một nhóm, số bạn thích bóng đá gấp đôi số bạn thích bóng chuyền. Có $6$ bạn thích cả hai môn và $30$ bạn thích ít nhất một trong hai môn. Hỏi có bao nhiêu bạn thích bóng đá, bao nhiêu bạn thích bóng chuyền?
\end{probbox}
\probsol{
Số bạn thích bóng đá $+$ số bạn thích bóng chuyền $=30+6=36$.\par Bóng chuyền $1$ phần, bóng đá $2$ phần: $3$ phần là $36$, một phần là $12$.\par
Bóng chuyền: $12$ bạn; bóng đá: $24$ bạn.
\dapso{Bóng đá $24$ bạn, bóng chuyền $12$ bạn.}}
\gap
\begin{probbox}{20}{\relax}
An, Bình và Cường có $60$ viên bi. An cho Bình $5$ viên, Bình cho Cường $8$ viên, rồi Cường cho An $3$ viên thì số bi của ba bạn bằng nhau. Hỏi lúc đầu mỗi bạn có bao nhiêu viên bi?
\end{probbox}
\probsol{
Lúc cuối mỗi bạn có $60:3=20$ (viên).\par
An: $20-3+5=22$ (viên).\par Bình: $20+8-5=23$ (viên).\par Cường: $20-8+3=15$ (viên).\par Thử lại: $22+23+15=60$.
\dapso{An $22$, Bình $23$, Cường $15$ viên bi.}}
\gap
\begin{probbox}{21}{\relax}
Hai hộp có $48$ chiếc bút. Chuyển $7$ chiếc bút từ hộp A sang hộp B thì số bút hai hộp bằng nhau. Hỏi lúc đầu mỗi hộp có bao nhiêu chiếc bút?
\end{probbox}
\probsol{
Lúc cuối mỗi hộp có $48:2=24$ (chiếc).\par Hộp A: $24+7=31$ (chiếc); hộp B: $24-7=17$ (chiếc).
\dapso{Hộp A $31$ chiếc, hộp B $17$ chiếc.}}
\gap
\begin{probbox}{22}{\relax}
Ba ngăn tủ có $90$ quyển sách. Chuyển $8$ quyển từ ngăn I sang ngăn II, rồi $12$ quyển từ ngăn II sang ngăn III, sau đó $5$ quyển từ ngăn III sang ngăn I thì số sách ba ngăn bằng nhau. Hỏi lúc đầu mỗi ngăn có bao nhiêu quyển sách?
\end{probbox}
\probsol{
Lúc cuối mỗi ngăn có $90:3=30$ (quyển).\par
Ngăn I: $30-5+8=33$ (quyển).\par Ngăn II: $30+12-8=34$ (quyển).\par Ngăn III: $30-12+5=23$ (quyển).
\dapso{$33$ quyển, $34$ quyển, $23$ quyển.}}
\gap
\begin{probbox}{23}{\relax}
Hà cho Lan $4$ hình dán, sau đó Lan cho Hà $10$ hình dán. Lúc này Hà có $26$ hình, Lan có $18$ hình. Hỏi lúc đầu mỗi bạn có bao nhiêu hình dán?
\end{probbox}
\probsol{
Trước khi Lan cho: Hà có $26-10=16$ (hình), Lan có $18+10=28$ (hình).\par
Trước khi Hà cho (lúc đầu): Hà có $16+4=20$ (hình), Lan có $28-4=24$ (hình).
\dapso{Hà $20$ hình, Lan $24$ hình.}}
\gap
\begin{probbox}{24}{\relax}
Mẹ chuyển từ rổ A sang rổ B số quả cam đúng bằng số cam rổ B đang có. Lúc này mỗi rổ có $30$ quả cam. Hỏi lúc đầu mỗi rổ có bao nhiêu quả cam?
\end{probbox}
\probsol{
Số cam rổ B đã gấp đôi, nên lúc đầu rổ B có $30:2=15$ (quả).\par Rổ A đã cho đi $15$ quả, nên lúc đầu có $30+15=45$ (quả).
\dapso{Rổ A $45$ quả, rổ B $15$ quả.}}
\gap
\begin{probbox}{25}{\relax}
Hai bể chứa $80$ lít nước. Bể A đổ sang bể B một lượng nước bằng lượng nước bể B đang có; sau đó bể B đổ sang bể A một lượng nước bằng lượng nước bể A lúc đó. Lúc này hai bể có lượng nước bằng nhau. Hỏi lúc đầu mỗi bể chứa bao nhiêu lít nước?
\end{probbox}
\probsol{
Lúc cuối mỗi bể có $40$ lít.\par
Trước lần đổ thứ hai: bể A có $40:2=20$ (lít), bể B có $40+20=60$ (lít).\par
Trước lần đổ thứ nhất: bể B có $60:2=30$ (lít), bể A có $20+30=50$ (lít).
\dapso{Bể A $50$ lít, bể B $30$ lít.}}
\gap
\begin{probbox}{26}{\relax}
Hùng có $36$ viên bi để trong ba hộp. Hùng chuyển từ hộp một sang hộp hai số bi đúng bằng số bi hộp hai đang có, rồi chuyển từ hộp hai sang hộp ba số bi đúng bằng số bi hộp ba đang có thì số bi ba hộp bằng nhau. Hỏi lúc đầu mỗi hộp có bao nhiêu viên bi?
\end{probbox}
\probsol{
Lúc cuối mỗi hộp có $36:3=12$ (viên).\par
Trước lần chuyển thứ hai: hộp ba có $12:2=6$ (viên), hộp hai có $12+6=18$ (viên), hộp một có $12$ viên.\par
Trước lần chuyển thứ nhất: hộp hai có $18:2=9$ (viên), hộp một có $12+9=21$ (viên).
\dapso{$21$ viên, $9$ viên, $6$ viên.}}
\gap
\begin{probbox}{27}{\relax}
Hòa, Bình và Việt có $70$ thẻ đọc. Hòa cho Bình $5$ thẻ, Bình cho Việt $7$ thẻ, rồi Việt cho Hòa $9$ thẻ thì lúc này Hòa và Bình có số thẻ bằng nhau và mỗi bạn có số thẻ gấp đôi số thẻ của Việt. Hỏi lúc đầu mỗi bạn có bao nhiêu thẻ đọc?
\end{probbox}
\probsol{
Lúc cuối: Việt $1$ phần, Hòa $2$ phần, Bình $2$ phần; $5$ phần là $70$ thẻ, một phần là $14$ thẻ.\par
Lúc cuối: Hòa $28$, Bình $28$, Việt $14$ thẻ.\par
Lúc đầu: Hòa $28-9+5=24$; Bình $28+7-5=30$; Việt $14-7+9=16$ (thẻ).
\dapso{Hòa $24$, Bình $30$, Việt $16$ thẻ.}}
\gap
\begin{probbox}{28}{\relax}
Một bác bán một nửa số xoài và thêm $2$ quả; sau đó bán một nửa số xoài còn lại và thêm $3$ quả thì còn lại $5$ quả. Hỏi lúc đầu bác có bao nhiêu quả xoài?
\end{probbox}
\probsol{
Trước lần bán thứ hai: $(5+3)\times2=16$ (quả).\par
Lúc đầu: $(16+2)\times2=36$ (quả).
\dapso{$36$ quả xoài.}}
\gap
\begin{probbox}{29}{\relax}
Ân và Bảo có $64$ viên bi. Ân cho Bảo số bi bằng số bi Bảo đang có; sau đó Bảo cho Ân số bi bằng số bi Ân đang có; rồi Ân lại cho Bảo số bi bằng số bi Bảo đang có. Lúc này số bi hai bạn bằng nhau. Hỏi lúc đầu mỗi bạn có bao nhiêu viên bi?
\end{probbox}
\probsol{
Lúc cuối mỗi bạn có $32$ viên.\par
Trước lần cho thứ ba: Bảo $16$, Ân $48$.\par
Trước lần cho thứ hai: Ân $24$, Bảo $40$.\par
Lúc đầu: Bảo $20$, Ân $44$.
\dapso{Ân $44$ viên, Bảo $20$ viên.}}
\gap
\begin{probbox}{30}{\relax}
Ngày thứ nhất dùng $\dfrac14$ lượng nước trong bể; ngày thứ hai dùng $\dfrac13$ lượng nước còn lại. Bể còn $60$ lít. Hỏi lúc đầu bể có bao nhiêu lít nước?
\end{probbox}
\probsol{
$60$ lít là $1-\dfrac13=\dfrac23$ lượng nước còn sau ngày thứ nhất: $60:\dfrac23=90$ (lít).\par
$90$ lít là $1-\dfrac14=\dfrac34$ lượng nước trong bể: $90:\dfrac34=120$ (lít).
\dapso{$120$ lít.}}
\dividerbar{Thử thách -- Nâng cao}
\begin{probbox}{31}{\stars{2}}
Trong các số từ $1$ đến $30$, có bao nhiêu số chia hết cho $2$, cho $3$ hoặc cho $5$? Có bao nhiêu số không chia hết cho số nào trong ba số đó?
\end{probbox}
\probsol{
Chia hết cho $2$: $15$; cho $3$: $10$; cho $5$: $6$.\par
Cho $6$: $5$; cho $10$: $3$; cho $15$: $2$; cho $30$ (cả ba): $1$.\par
Điền từ trong ra ngoài:
\begin{center}
\stpVennIII{cho $2$}{cho $3$}{cho $5$}{$8$}{$4$}{$2$}{$4$}{$2$}{$1$}{$1$}{không: $8$}
\end{center}
Chia hết cho ít nhất một số: $15+10+6-5-3-2+1=22$.\par Không chia hết cho số nào: $30-22=8$.
\dapso{$22$ số; $8$ số.}}
\gap
\begin{probbox}{32}{\stars{1}}
Một lớp có $40$ học sinh, $25$ bạn thích Toán và $28$ bạn thích Văn. Số bạn thích cả hai môn ít nhất là bao nhiêu? Nhiều nhất là bao nhiêu?
\end{probbox}
\probsol{
Số bạn thích ít nhất một môn không quá $40$, nên số bạn thích cả hai ít nhất là $25+28-40=13$ (khi bạn nào cũng thích ít nhất một môn).\par
Số bạn thích cả hai không thể nhiều hơn nhóm nhỏ hơn: nhiều nhất là $25$ (khi mọi bạn thích Toán đều thích Văn).
\dapso{Ít nhất $13$ bạn, nhiều nhất $25$ bạn.}}
\gap
\begin{probbox}{33}{\stars{1}}
$45$ học sinh, mỗi bạn tham gia ít nhất một trong ba câu lạc bộ A, B, C. CLB A có $24$ bạn, B có $22$ bạn, C có $20$ bạn; $9$ bạn tham gia A và B, $7$ bạn tham gia B và C, $8$ bạn tham gia C và A. Hỏi có bao nhiêu bạn tham gia cả ba câu lạc bộ?
\end{probbox}
\probsol{
$45=24+22+20-9-7-8+(\text{cả ba})=42+(\text{cả ba})$.\par
Cả ba câu lạc bộ: $45-42=3$ (bạn).
\dapso{$3$ bạn.}}
\gap
\begin{probbox}{34}{\stars{2}}
Trong đội tuyển của một trường có $30$ bạn thi Toán, $24$ bạn thi Tiếng Anh và $20$ bạn thi Khoa học. Biết có $46$ bạn thi Toán hoặc Tiếng Anh, $38$ bạn thi Tiếng Anh hoặc Khoa học, $42$ bạn thi Khoa học hoặc Toán và không có bạn nào thi cả ba môn.\par
\pa{a} Có bao nhiêu bạn thi cả Toán và Tiếng Anh? \quad \pa{b} cả Tiếng Anh và Khoa học? \quad \pa{c} cả Khoa học và Toán?\par
\pa{d} Đội tuyển có bao nhiêu học sinh?
\end{probbox}
\probsol{
\pa{a} $30+24-46=8$ (bạn).\par
\pa{b} $24+20-38=6$ (bạn).\par
\pa{c} $20+30-42=8$ (bạn).\par
\pa{d} $30+24+20-8-6-8=52$ (học sinh).
\dapso{\pa{a} $8$. \pa{b} $6$. \pa{c} $8$. \pa{d} $52$ học sinh.}}
\gap
\begin{probbox}{35}{\stars{1}}
Trong một lớp có $8$ bạn thích cả chó và mèo, $5$ bạn không thích con nào, $14$ bạn thích chó. Số bạn chỉ thích mèo gấp đôi số bạn chỉ thích chó. Hỏi lớp có bao nhiêu học sinh và có bao nhiêu bạn thích mèo?
\end{probbox}
\probsol{
Chỉ thích chó: $14-8=6$ (bạn).\par Chỉ thích mèo: $6\times2=12$ (bạn).\par
Số học sinh của lớp: $6+12+8+5=31$ (học sinh).\par Thích mèo: $12+8=20$ (bạn).
\dapso{$31$ học sinh; $20$ bạn thích mèo.}}
\gap
\begin{probbox}{36}{\stars{1}}
Trong các số từ $1$ đến $100$, có bao nhiêu số chia hết cho $3$ nhưng không chia hết cho $4$?
\end{probbox}
\probsol{
Chia hết cho $3$: $33$ số (vì $100=3\times33+1$).\par Chia hết cho cả $3$ và $4$, tức là chia hết cho $12$: $8$ số ($12;\ 24;\ \dots;\ 96$).\par
Chia hết cho $3$ nhưng không chia hết cho $4$: $33-8=25$ (số).
\dapso{$25$ số.}}
\gap
\begin{probbox}{37}{\stars{2}}
Mỗi bạn trong một nhóm chơi ít nhất một trong ba môn thể thao A, B, C: $55$ bạn chơi A, $45$ bạn chơi B, $40$ bạn chơi C. Có $20$ bạn chơi đúng hai môn và $5$ bạn chơi cả ba môn. Hỏi nhóm có bao nhiêu bạn?
\end{probbox}
\probsol{
Trong tổng $55+45+40=140$, bạn chơi một môn được tính một lần, bạn chơi đúng hai môn được tính hai lần, bạn chơi ba môn được tính ba lần.\par
Số bạn chơi đúng một môn: $140-2\times20-3\times5=85$ (bạn).\par
Số bạn của nhóm: $85+20+5=110$ (bạn).
\dapso{$110$ bạn.}}
\gap
\begin{probbox}{38}{\stars{1}}
Trong $50$ học sinh, có $35$ bạn biết bơi và $30$ bạn biết chơi cờ vua.\par
\pa{a} Có ít nhất bao nhiêu bạn biết cả hai? \quad \pa{b} Nếu có đúng $4$ bạn không biết môn nào thì có bao nhiêu bạn biết cả hai?
\end{probbox}
\probsol{
\pa{a} Số bạn biết ít nhất một môn không quá $50$, nên số bạn biết cả hai ít nhất là $35+30-50=15$.\par
\pa{b} Biết ít nhất một môn: $50-4=46$ (bạn); biết cả hai: $35+30-46=19$ (bạn).
\dapso{\pa{a} $15$ bạn. \pa{b} $19$ bạn.}}
\gap
\begin{probbox}{39}{\stars{2}}
Quân, Sơn và Tùng có $48$ viên bi. Quân lấy một nửa số bi của mình chia đều cho hai bạn còn lại; sau đó Sơn cũng lấy một nửa số bi mình đang có chia đều cho hai bạn còn lại; cuối cùng Tùng cũng làm như vậy. Lúc này số bi của ba bạn bằng nhau. Hỏi lúc đầu mỗi bạn có bao nhiêu viên bi?
\end{probbox}
\probsol{
Lúc cuối mỗi bạn có $16$ viên. Người cho giữ lại một nửa; mỗi bạn kia nhận một phần tư số bi người cho có trước đó.
\begin{center}
\renewcommand{\arraystretch}{1.25}
\begin{tabular}{|l|c|c|c|}
\hline
 & Quân & Sơn & Tùng\\\hline
Lúc cuối & $16$ & $16$ & $16$\\\hline
Trước khi Tùng chia & $16-8=8$ & $16-8=8$ & $16\times2=32$\\\hline
Trước khi Sơn chia & $8-4=4$ & $8\times2=16$ & $32-4=28$\\\hline
Lúc đầu & $4\times2=\mathbf{8}$ & $16-2=\mathbf{14}$ & $28-2=\mathbf{26}$\\\hline
\end{tabular}
\end{center}
\dapso{Quân $8$, Sơn $14$, Tùng $26$ viên bi.}}
\gap
\begin{probbox}{40}{\stars{2}}
Có ba thùng dầu A, B, C. Người ta chia số dầu thùng A thành $3$ phần bằng nhau, rót sang thùng B một phần, thùng C một phần. Sau đó chia số dầu thùng B lúc đó thành $4$ phần bằng nhau, rót sang thùng A một phần, thùng C một phần. Rồi chia số dầu thùng C lúc đó thành $3$ phần bằng nhau, rót sang thùng A một phần, thùng B một phần. Lúc này thùng A có $40$ lít, thùng B có $46$ lít và thùng C có $18$ lít. Hỏi lúc đầu mỗi thùng có bao nhiêu lít dầu?
\end{probbox}
\probsol{
Trước khi thùng C rót: C còn lại $1$ phần trong $3$ phần, nên C có $18\times3=54$ (lít); A và B mỗi thùng nhận $18$ lít: A có $40-18=22$ (lít), B có $46-18=28$ (lít).\par
Trước khi thùng B rót: B còn lại $2$ phần trong $4$ phần, nên B có $28\times2=56$ (lít); A và C mỗi thùng nhận $56:4=14$ (lít): A có $22-14=8$ (lít), C có $54-14=40$ (lít).\par
Trước khi thùng A rót: A còn lại $1$ phần trong $3$ phần, nên A có $8\times3=24$ (lít); B và C mỗi thùng nhận $8$ lít: B có $56-8=48$ (lít), C có $40-8=32$ (lít).
\dapso{Thùng A $24$ lít, thùng B $48$ lít, thùng C $32$ lít.}}
\gap
\begin{probbox}{41}{\stars{1}}
Lan đọc một quyển sách. Ngày thứ nhất Lan đọc $\dfrac13$ quyển sách và thêm $4$ trang; ngày thứ hai đọc một nửa số trang còn lại và thêm $6$ trang; ngày thứ ba đọc nốt $24$ trang cuối. Hỏi quyển sách có bao nhiêu trang?
\end{probbox}
\probsol{
Ngày thứ hai đọc một nửa số trang còn lại và thêm $6$ trang thì còn $24$ trang, nên trước ngày thứ hai còn $(24+6)\times2=60$ (trang).\par
Ngày thứ nhất: $60+4=64$ (trang) là $1-\dfrac13=\dfrac23$ quyển sách.\par Quyển sách có: $64:\dfrac23=96$ (trang).
\dapso{$96$ trang.}}
\gap
\begin{probbox}{42}{\stars{2}}
Ba hộp I, II, III đựng kẹo. Đầu tiên hộp I cho hộp II và hộp III mỗi hộp số kẹo đúng bằng số kẹo hộp đó đang có. Sau đó hộp II cho hộp I và hộp III mỗi hộp số kẹo đúng bằng số kẹo hộp đó lúc ấy. Lúc này hộp I có $20$ cái, hộp II có $30$ cái, hộp III có $40$ cái kẹo. Hỏi lúc đầu mỗi hộp có bao nhiêu cái kẹo?
\end{probbox}
\probsol{
Trước khi hộp II cho: hộp I có $20:2=10$, hộp III có $40:2=20$, hộp II có $30+10+20=60$ (cái).\par
Trước khi hộp I cho: hộp II có $60:2=30$, hộp III có $20:2=10$, hộp I có $10+30+10=50$ (cái).\par
Thử lại: $(50;\ 30;\ 10)\to(10;\ 60;\ 20)\to(20;\ 30;\ 40)$.
\dapso{Hộp I $50$ cái, hộp II $30$ cái, hộp III $10$ cái.}}
\gap
\begin{probbox}{43}{\stars{1}}
Bình dùng $\dfrac14$ số tiền để mua sách, rồi dùng $\dfrac25$ số tiền còn lại để mua cặp, sau đó tiêu $15\,000$ đồng ăn sáng. Bình còn lại $30\,000$ đồng. Hỏi lúc đầu Bình có bao nhiêu tiền?
\end{probbox}
\probsol{
Trước khi ăn sáng: $30\,000+15\,000=45\,000$ (đồng).\par
Số tiền này là $1-\dfrac25=\dfrac35$ số tiền còn lại sau khi mua sách: $45\,000:\dfrac35=75\,000$ (đồng).\par
Số tiền đó là $\dfrac34$ số tiền lúc đầu: $75\,000:\dfrac34=100\,000$ (đồng).
\dapso{$100\,000$ đồng.}}
\gap
\begin{probbox}{44}{\stars{1}}
Mai nói: ``Lấy tuổi của tớ nhân với $3$, cộng $12$, chia cho $5$ rồi trừ đi $4$ thì được $5$.'' Hỏi Mai bao nhiêu tuổi?
\end{probbox}
\probsol{
$5+4=9$; \ $9\times5=45$; \ $45-12=33$; \ $33:3=11$.\par
Thử lại: $11\times3+12=45$; $45:5=9$; $9-4=5$.
\dapso{$11$ tuổi.}}
\gap
\begin{probbox}{45}{\stars{2}}
$40$ học sinh, mỗi bạn tham gia ít nhất một trong ba câu lạc bộ: Đọc sách $22$ bạn, Khoa học $18$ bạn, Múa $20$ bạn. Có $3$ bạn tham gia cả ba câu lạc bộ. Hỏi có bao nhiêu bạn tham gia đúng hai câu lạc bộ? Đúng một câu lạc bộ?
\end{probbox}
\probsol{
Trong tổng $22+18+20=60$: bạn tham gia một CLB được tính một lần, hai CLB được tính hai lần, ba CLB được tính ba lần.\par
Bỏ ra $3$ bạn tham gia cả ba CLB: còn $37$ bạn, được tính $60-9=51$ lần.\par
Mỗi bạn tham gia đúng hai CLB được tính thừa $1$ lần: số bạn tham gia đúng hai CLB là $51-37=14$.\par Đúng một CLB: $37-14=23$ (bạn).
\dapso{Đúng hai CLB: $14$ bạn; đúng một CLB: $23$ bạn.}}
\gap
\begin{probbox}{46}{\stars{2}}
Trong một lớp, $\dfrac35$ số học sinh thích Toán, $\dfrac12$ số học sinh thích Tiếng Anh, $\dfrac14$ số học sinh thích cả hai môn và có $6$ bạn không thích môn nào trong hai môn đó. Hỏi lớp có bao nhiêu học sinh?
\end{probbox}
\probsol{
Thích ít nhất một môn: $\dfrac35+\dfrac12-\dfrac14=\dfrac{17}{20}$ số học sinh.\par
Không thích môn nào: $1-\dfrac{17}{20}=\dfrac{3}{20}$ số học sinh, ứng với $6$ bạn.\par
Số học sinh của lớp: $6:\dfrac{3}{20}=40$ (học sinh).
\dapso{$40$ học sinh.}}
\gap
\begin{probbox}{47}{\stars{1}}
Một cửa hàng buổi sáng bán $40\%$ số gạo, buổi chiều bán $50\%$ số gạo còn lại thì còn $66$ kg. Hỏi lúc đầu cửa hàng có bao nhiêu ki-lô-gam gạo?
\end{probbox}
\probsol{
Sau buổi sáng còn $60\%$ số gạo; buổi chiều bán một nửa số đó nên còn $30\%$ số gạo.\par
$30\%$ số gạo là $66$ kg, nên số gạo là: $66:30\times100=220$ (kg).
\dapso{$220$ kg.}}
\gap
\begin{probbox}{48}{\stars{1}}
An, Bình, Chi và Dũng có $100$ tấm thẻ. An cho Bình $3$ thẻ, Bình cho Chi $5$ thẻ, Chi cho Dũng $7$ thẻ, rồi Dũng cho An $9$ thẻ thì số thẻ của bốn bạn bằng nhau. Hỏi lúc đầu mỗi bạn có bao nhiêu tấm thẻ?
\end{probbox}
\probsol{
Lúc cuối mỗi bạn có $100:4=25$ (thẻ).\par
An: $25-9+3=19$; \ Bình: $25+5-3=27$; \ Chi: $25+7-5=27$; \ Dũng: $25-7+9=27$ (thẻ).
\dapso{An $19$, Bình $27$, Chi $27$, Dũng $27$ thẻ.}}
\gap
\begin{probbox}{49}{\stars{1}}
Có bao nhiêu số có hai chữ số mà có ít nhất một chữ số $5$?
\end{probbox}
\probsol{
Chữ số hàng chục là $5$: từ $50$ đến $59$, có $10$ số.\par Chữ số hàng đơn vị là $5$: $15;\ 25;\ \dots;\ 95$, có $9$ số.\par Cả hai chữ số là $5$: chỉ có số $55$.\par
Có ít nhất một chữ số $5$: $10+9-1=18$ (số).
\dapso{$18$ số.}}
\gap
\begin{probbox}{50}{\stars{2}}
An, Bình và Chi chơi ba ván. Sau mỗi ván, người thua phải cho mỗi người còn lại số bi đúng bằng số bi người đó đang có. An thua ván 1, Bình thua ván 2, Chi thua ván 3. Lúc cuối An có $16$ viên, Bình có $24$ viên và Chi có $32$ viên bi. Hỏi lúc đầu mỗi bạn có bao nhiêu viên bi?
\end{probbox}
\probsol{
Khi tính ngược, số bi của hai người thắng được chia đôi, người thua nhận lại phần đã cho.
\begin{center}
\renewcommand{\arraystretch}{1.25}
\begin{tabular}{|l|c|c|c|}
\hline
 & An & Bình & Chi\\\hline
Lúc cuối & $16$ & $24$ & $32$\\\hline
Trước ván 3 & $8$ & $12$ & $32+8+12=52$\\\hline
Trước ván 2 & $4$ & $12+4+26=42$ & $26$\\\hline
Lúc đầu & $4+21+13=\mathbf{38}$ & $\mathbf{21}$ & $\mathbf{13}$\\\hline
\end{tabular}
\end{center}
\dapso{An $38$, Bình $21$, Chi $13$ viên bi.}}
